\documentclass[14pt,a4paper]{extarticle}
\usepackage[T2A]{fontenc}
\usepackage[utf8]{inputenc}
\usepackage[english,russian]{babel}
\usepackage{lmodern}
\usepackage{geometry}
\usepackage{array}
\usepackage{tabularx}
\usepackage{longtable}
\usepackage{ltablex}
\usepackage{booktabs}
\usepackage{hyperref}
\usepackage{indentfirst}
\usepackage{microtype}
\usepackage{graphicx}
\usepackage{setspace}
\usepackage{listings}
\usepackage{graphicx}
\usepackage{caption}
\usepackage{float}
\usepackage{tocloft}
\usepackage{titlesec}
\usepackage{newtxtext,newtxmath}

\titleformat{\section}        % форматируем \section
  {\normalfont\bfseries}      % стиль: обычный + жирный
  {\thesection.}              % номер секции с точкой
  {2mm}                       % расстояние между номером и текстом
  {}                           % перед заголовком можно добавить что-то
\titlespacing*{\section}{1.5cm}{*3}{*1}  % {отступ слева}{отступ сверху}{отступ снизу}

\titleformat{\subsection}        % форматируем \section
  {\normalfont\bfseries}      % стиль: обычный + жирный
  {\thesubsection}              % номер секции с точкой
  {2mm}                       % расстояние между номером и текстом
  {}                           % перед заголовком можно добавить что-то
\titlespacing*{\subsection}{1.5cm}{*3}{*1}  % {отступ слева}{отступ сверху}{отступ снизу}

\captionsetup{justification=centering}
\renewcommand{\ttdefault}{cmtt}
\keepXColumns
\setlength{\parindent}{3em}
\newgeometry{left=2.5cm,right=2.5cm,top=3cm,bottom=2.5cm}

\renewcommand{\contentsname}{\hspace*{4em}{\smallСодержание}}
\setlength{\cftbeforesecskip}{0.5em}
\renewcommand{\cftsecleader}{\cftdotfill{\cftdotsep}}
\renewcommand{\cftsecaftersnum}{.\hspace{2mm}}

\title{Отчёт по "Программно-аппаратные средства интеллектуальных систем"}
\author{Выполнил студент группы М-РИТ-25-1, Станиславчук Сергей}

\begin{document}

\setstretch{1.5}
\noindent
\begin{minipage}[t]{0.01\textwidth}
    \includegraphics{~/programmer/++Programmer/5_1/LSTU.jpg}
\end{minipage}%
\begin{minipage}[c]{1.1\linewidth}
    \centering
    {\small
    МИНИСТЕРСТВО НАУКИ И ВЫСШЕГО ОБРАЗОВАНИЯ РФ \\[1em]
    ФЕДЕРАЛЬНОЕ ГОСУДАРСТВЕННОЕ БЮДЖЕТНОЕ \\
    ОБРАЗОВАТЕЛЬНОЕ УЧРЕЖДЕНИЕ ВЫСШЕГО ОБРАЗОВАНИЯ \\
    «ЛИПЕЦКИЙ ГОСУДАРСТВЕННЫЙ ТЕХНИЧЕСКИЙ УНИВЕРСИТЕТ»
    }
\end{minipage}%

\vspace{3em}

\noindent
\hspace{1em}%
\begin{minipage}[t]{0.8\textwidth}
\textbf{Институт} \hspace{2em} компьютерных наук \\
\textbf{Кафедра} \hspace{2em} автоматизированных систем управления
\end{minipage}


\vspace{6em}

\begin{center}
\textbf{ИНДИВИДУАЛЬНОЕ ЗАДАНИЕ} \\
по дисциплине \\ <<Программно-аппаратные средства интеллектуальных систем>>
\end{center}

\vspace{6em}

\noindent
\begin{tabular*}{\textwidth}{@{\extracolsep{\fill}} l c r }
    Студент М-РИТ-25-1 & \makebox[4cm]{\hrulefill} & Станиславчук С.М. \\[2em]
Руководитель & \\
Ст. преподаватель & \makebox[4cm]{\hrulefill} & Болдырихин О.В. \\
\end{tabular*}

\vfill
\begin{center}
    Липецк, 2025
\end{center}

\thispagestyle{empty}

\newpage
\tableofcontents
\newpage

\section{Цель работы}
Исследование определенной архитектуры суперкомпьютерной системы из рейтинга TOP500.

\section{Задание}

\begin{sloppypar}
    \renewcommand{\labelitemi}{$-$}

    Задание 2. Описать архитектуру определенной суперкомпьютерной системы из рейтинга TOP500:

    \begin{itemize}
        \item основные характеристики: архитектура, разработчик, место инсталляции,
            процессор, сопроцессор, размер памяти, интерконнект;
        \item основные технико-эксплуатационные показатели;
        \item процессор: модель, характеристики, описание архитектуры;
        \item сопроцессор: модель, характеристики, описание архитектуры;
        \item интерконнект: хаарктеристики, описание;
        \item используемое программное обеспечение;
        \item область использования, решаемые задачи.
    \end{itemize}

    Рассчитать дополнительные показатели: эффективность архитектуры по различным тестам и др.
    Проанализировать архитектуру суперкомпьютерной системы на основе рассчитаныых значений.
\end{sloppypar}

\section{Описание суперкомпьютерной системы}

Выбранная суперкомпьютерная система: El Capitan (Hewlett Packard Enterprise El Capitan).

\subsection{Основные характеристики}

    Архитектура -- Cray EX255a~\cite{top500}.

    Разработчик -- HPE Cray, LLNL, AMD.

    Место инсталляции -- Ливерморская национальная лаборатория им. Лоуренса, Калифорния, США~\cite{wikiped}.

    Процессор -- AMD EPYC 4th Gen (Zen4) 24C «Genoa», 1{,}8 ГГц ($1\,051\,392$ ядра).

    Сопроцессор -- AMD Instinct MI300A, $43\,808$ APU ($9\,988\,224$ вычислительных блоков).

    Размер памяти -- $5{,}4375$ ПБ (128 ГБ на сокет).

    Интерконнект -- HPE Cray Slingshot 11.

\subsection{Основные технико-эксплуатационные показатели}

    Rmax = $1\,809.00$ ПФлоп/с

    Power = $29\,684.6$ кВт

    Rpeak = $2\,821\,104\,000$ ПФлоп/с

    Nmax = $24\,768\,000$

    HPCG = $17\,406.00$ ТФлоп/с

\subsection{Подробное описание характеристик}

    Процессор(ы) -- AMD 4th Gen (Zen4) x86 EPYC 24C <<Genoa>> с тактовой частотой 1.8 ГГц ($1\,051\,392$ ядра).
    Поддерживает X-86 инструкции и AVX-512 и BFloat16 расширения. Каждый ССD
    (ПЗС, прибор с зарядовой связью) имеет 8 Zen ядер которые совместно
    используют 32 мб L3 кэш.

    Сопроцессор(ы) -- $43\,808$ графических процессоров AMD Instinct MI300A
    ($9\,988\,224$ вычислительных блока (по 228 в каждом)). В каждом узле содержиться 4 AMD MI300A APU.
    Каждый APU имеет 6 XCD (Accelerator Complex Die), которые служат как GPU.
    AMD MI300A R-Peak - 61.3 терафлопс в векторной системе двойной точности.

    В общей сложности $649\,246\,336$ потоковых процессоров (по $14\,592$ в каждом)

    Размер памяти -- $5.4375$ петабайт (128 гб на сокет)

    Интерконнет -- Infinity Fabric (IF). Коммутаторы HPE Slingshot Interconnect - используют топологию Dragonfly.

    Широкополосный коммутатор с 64 портами и двунаправленный пропускной
    способностью 25.6 Тб/с. Каждый порт может обеспеченивать однонаправленную
    пропускную способность до 200 Гбит/с (25 Гб/с). Четыре интерфейса HPE Slingshot 200 GbE на каждый узел. Общая пропускная способность ввода данных в узлы составляет 100 Гб/с~\cite{hardware}

    Используемое ПО -- Tri-Lab Operating System Software (TOSS), система,
    основанная на линукс для NNSA лабораторий, для проведения сложных
    физических симуляций для управления ядерными запасами Соединенных Штатов.

    \subsection{Эффективность архитектуры по различным тестам}

    Энергоэффективность (Гфлоп/Вт) = 60.94 (из рейтинга)

    Эффективность по HPCG: $E_{ HPCG } = \frac{HPCG}{R_{peak}} = \frac{17.4}{2821} = 0.0062 = 0.62\%$

    Эффективность по LINPACK: $E_{ HPL } = \frac{R_{max}}{R_{peak}} * 100\% = \frac{1809}{2821} = 0.64 = 64\%$

\newpage
\section{Анализ архитектуры суперкомпьютерной системы}

Высокое значение эффективности по тесту LINPACK (64\% от $R_{peak}$ в
значительной степени обусловлено архитектурными особенностями графических
ускорителей AMD Instinct MI300A. Данные ускорители обеспечивают высокую
производительность операций двойной точности и ориентированы на выполнение
плотных численных вычислений с высокой арифметической интенсивностью. Алгоритм
HPL характеризуется регулярным доступом к памяти и большим числом операций с
плавающей точкой на один элемент данных, что позволяет эффективно загружать
вычислительные блоки GPU и минимизировать простои потоковых процессоров.

В то же время низкая эффективность по тесту HPCG (0.62\% от $R_{ peak }$)
объясняется тем, что данный бенчмарк основан на разреженных матрицах и
нерегулярных шаблонах доступа к памяти. Такие нагрузки плохо масштабируются на
GPU-архитектурах, поскольку приводят к дивергенции потоков и снижению
вычислительной загрузки, вследствие чего большая часть вычислительных ресурсов
остается недоиспользованной.

Роль CPU. Процессоры AMD EPYC архитектуры Zen 4 в системе El Capitan
функционируют преимущественно как управляющие и координирующие элементы, а не
как основные вычислители. Низкая тактовая частота (1.8 ГГц) компенсируется
поддержкой современных векторных расширений (AVX-512, BFloat16), однако в
условиях HPCG-подобных нагрузок, чувствительных к задержкам и
производительности одного потока, это ограничивает вклад CPU в общее ускорение
вычислений. В результате центральные процессоры не способны компенсировать
снижение эффективности GPU при выполнении нерегулярных и память-зависимых
операций.

\section{Вывод}

В ходе выполненного индивидуального домашнего задания была изучена и
проанализована система самого мощного суперкомпьюетера (на момент ноябрь 2025
года) El Capitan.

\newpage

\begin{thebibliography}{9}

    \bibitem{top500}
        HPCG - NOVEMBER 2025
        \url{https://top500.org/lists/hpcg/2025/11/}

    \bibitem{wikiped}
        El Capitan (supercomputer)
        \url{https://en.wikipedia.org/wiki/El_Capitan_(supercomputer)}

    \bibitem{hardware}
        Using El Capitan Systems: Hardware Overview
        \url{https://hpc.llnl.gov/documentation/user-guides/using-el-capitan-systems/hardware-overview}
\end{thebibliography}

\end{document}
